\documentclass[aspectratio=169]{beamer}

\usetheme{Madrid}
\usecolortheme{default}
\setbeamertemplate{navigation symbols}{}

\usepackage{amsmath, amssymb, bm}
\usepackage{graphicx}

% ---------- Macros ----------
\newcommand{\R}{\mathbb{R}}
\newcommand{\uu}{\bm{u}}
\newcommand{\pp}{p}
\newcommand{\Rey}{\mathrm{Re}}
\newcommand{\nuu}{\nu}
\newcommand{\tmark}[1]{\hfill{\scriptsize\textcolor{gray}{#1}}}

% Uncomment if you want speaker notes:
% \usepackage{pgfpages}
% \setbeameroption{show notes on second screen=right}

\title[Turbulence: What is it?]{What is Turbulence?\\Key Features, Examples, and Why It Matters}
\author{
\href{https://cassyni.com/events/HsafxyMiJvA1JVCU9eef3T/seminar}{(adapted from Steve Brunton lecture)}}
\date{}

\begin{document}
% -------------------------
\begin{frame}
  \titlepage
\end{frame}

% -------------------------
\begin{frame}{Lecture roadmap}
\begin{itemize}
  \item Why turbulence is hard (and why it matters)
  \item Canonical flows: a ``vocabulary'' for turbulence
  \item Examples: pipe, wake, boundary layer, cavity, jets, mixing layers
  \item ``More complex'' flows (geometry, separation, multi-scale physics)
  \item Open-source resources: databases and CFD codes
\end{itemize}
\end{frame}


% -------------------------
\begin{frame}{Why turbulence is hard}
\begin{columns}[T,onlytextwidth]
\column{0.56\textwidth}
\begin{itemize}
  \item Turbulence is a \emph{multi-scale}, \emph{history-dependent} dynamics problem.
  \item Small perturbations can grow, interact, and reorganize the flow.
  \item Many flows are \emph{linearly stable} yet become turbulent through finite-amplitude disturbances.
  \item Predicting transition and statistics (drag, mixing, heat transfer, intermittency) is still challenging.
\end{itemize}

\column{0.44\textwidth}
\includegraphics[width=\textwidth]{pictures/5-1/turbulence.png}

\vspace{0.35em}
{\footnotesize A classic motivation appears in the \emph{Feynman Lectures} (Vol.\ II) discusses why predicting real flows is so difficult.}
\end{columns}
\end{frame}

% -------------------------
\begin{frame}{Canonical flows (Van Dyke ``album'')}
\begin{columns}[T,onlytextwidth]
\column{0.55\textwidth}
\begin{block}{Idea}
A small set of ``canonical'' configurations show up \emph{everywhere}.  Learn these and you gain a language for:
\begin{itemize}
  \item transition mechanisms (instability \(\rightarrow\) breakdown),
  \item coherent structures (vortices, streaks, shear layers),
  \item scaling laws and nondimensional groups.
\end{itemize}
\end{block}
\column{0.4\textwidth}
\includegraphics[width=0.75\textwidth]{pictures/5-1/album.png}
\end{columns}
\end{frame}

% -------------------------
\begin{frame}{A minimal tool: Reynolds number}
\begin{itemize}
  \item Reynolds number compares inertia to viscosity:
  \[
    \mathrm{Re}=\frac{UL}{\nu}.
  \]
  \item ``High'' \(\mathrm{Re}\) often means stronger instabilities, thinner boundary layers, and wider separation of scales.
  \item But \(\mathrm{Re}\) alone does \emph{not} determine the outcome:
  geometry, forcing, roughness, and disturbance history matter.
\end{itemize}
\end{frame}

% =========================
% Canonical flows
% =========================

\section{Canonical flows}

% -------------------------
\begin{frame}{Pipe flow}
\begin{columns}[T,onlytextwidth]
\column{0.5\textwidth}
\begin{itemize}
  \item Classic internal flow: long straight pipe, pressure-driven.
  \item Laminar solution exists for all \(\mathrm{Re}\), but becomes \emph{fragile} to disturbances.
  \item Transition often involves localized ``puffs'' and ``slugs'' (intermittency).
  \item Active control is possible: carefully designed perturbations can suppress turbulence (drag reduction).
\end{itemize}
\column{0.45\textwidth}
\includegraphics[width=\textwidth]{pictures/5-1/t_0300.jpg}
\end{columns}
\end{frame}

%-------------------------------
\begin{frame}{Pipe flow: Slugs and puffs}
\begin{itemize}
\item At high velocities, the flow is turbulent, and at intermediate velocities, transitional. 
\item The transitional flow is a mix of eddying flows and laminar flows. 
\item The eddying flows come in different varieties. 
\begin{itemize}
\item "Slugs" grow aggressively as they flow downstream; 
\item "Puffs" maintain a fixed size as they flow downstream. 
\end{itemize}
\end{itemize}

\centering
\includegraphics[width=0.7\textwidth]{pictures/5-1/turbulence-pipe.png}

\vspace{0.35em}
{\footnotesize Source: Graham, M. Turbulence spreads like wildfire. Nature 526, 508–509 (2015).}
\end{frame}

%-------------------
\begin{frame}{Puffs and slugs}
\begin{itemize}
\item Schematic of energy spectra corresponding to the flow inside the marked regions. 
\item Regardless of the flow variety, the energy spectrum is universal for small eddies.
\end{itemize}
\centering
\includegraphics[width=0.75\textwidth]{pictures/5-1/1-goingwiththe.jpg}

\vspace{0.35em}
{\footnotesize Source: https://phys.org/news/2020-01-insights-mysterious-fluid-motions.html}

\end{frame}

% -------------------------
\begin{frame}{Wake flows}
\begin{columns}[T,onlytextwidth]
\column{0.5\textwidth}
\begin{itemize}
  \item Flow past a body leaves a \emph{wake} with a velocity deficit and strong shear.
  \item Vortex shedding (\(\rightarrow\) K\'arm\'an vortex street) produces unsteady loading and noise.
  \item Wakes are central to:
  \begin{itemize}
    \item aircraft wing tips and propellers,
    \item wind farms and arrays,
    \item animal locomotion and swimming.
  \end{itemize}
\end{itemize}
\column{0.45\textwidth}
\includegraphics[width=\textwidth]{pictures/5-1/t_0360.jpg}
\end{columns}
\end{frame}


% -------------------------
\begin{frame}{Boundary layer}
\begin{columns}[T,onlytextwidth]
\column{0.56\textwidth}
\begin{itemize}
  \item Near a wall, no-slip forces a thin region of strong shear: the boundary layer.
  \item Key questions:
  \begin{itemize}
    \item when does laminar flow transition to turbulence?
    \item how do near-wall streaks and vortices sustain turbulence?
    \item how do roughness and pressure gradients change drag?
  \end{itemize}
  \item Boundary layers dominate drag and heat transfer on vehicles, buildings, and animals.
\end{itemize}
\column{0.44\textwidth}
\includegraphics[width=\textwidth]{pictures/5-1/turbulent-boundary.png}
{\footnotesize Source: Dou, Hua-Shu. (2018). Singularity of Navier-Stokes Equations. 10.48550/arXiv.1805.12053.} 
\end{columns}
\end{frame}

% -------------------------
\begin{frame}{Boundary layer: spatial development}
\centering
\includegraphics[width=0.92\textwidth]{pictures/5-1/boundary-layer.png}
\vspace{0.35em}

{\footnotesize Source: https://www.grasacoustics.com/blog/aerodynamic-measurements-components-of-the-boundary-layer}
\end{frame}

% -------------------------
\begin{frame}{Cavity flow}
\begin{columns}[T,onlytextwidth]
\column{0.58\textwidth}
\begin{itemize}
  \item A shear layer forms over a cavity opening.
  \item Feedback between shear-layer instability and acoustic/pressure waves can create resonant oscillations.
  \item Applications:
  \begin{itemize}
    \item flow over developing heart chambers, aneurysms,
    \item gaps, steps, and surface defects,
    \item flow-induced vibration.
  \end{itemize}
\end{itemize}
\column{0.42\textwidth}
\includegraphics[width=\textwidth]{pictures/5-1/aneurysm.png}
{\footnotesize Source: Boniforti MA, Magini R, Orosco Salinas T. Hemodynamic Investigation of the Flow Diverter Treatment of Intracranial Aneurysm. Fluids. 2023; 8(7):189. }
\end{columns}
\end{frame}



% =========================
% More complex flows and resources
% =========================


% -------------------------
\begin{frame}{Open data: turbulence databases}
\begin{columns}[T,onlytextwidth]
\column{0.57\textwidth}
\begin{itemize}
  \item Public datasets let you test ideas without running a massive DNS yourself.
  \item \href{https://turbulence.idies.jhu.edu/home}{Example: Johns Hopkins Turbulence Databases (JHTDB)}  
  \begin{itemize}
    \item isotropic turbulence, channel flow, boundary layers, geophysical, \ldots
    \item query-able via web/API-style interfaces
    \item great for teaching statistics, structure functions, spectra, and coherent-structure detection
  \end{itemize}
\end{itemize}
\column{0.43\textwidth}
\includegraphics[width=\textwidth]{pictures/5-1/JHU-dataset.png}
\end{columns}
\end{frame}

% -------------------------
\begin{frame}{Open-source CFD codes}
\begin{columns}[T,onlytextwidth]
\column{0.52\textwidth}
\begin{block}{Two common ecosystems}
\begin{itemize}
  \item \href{https://www.openfoam.com/}{\textbf{OpenFOAM}}: finite-volume framework with many RANS/LES solvers; strong community support.
  \item \href{https://nek5000.mcs.anl.gov/}{\textbf{Nek5000}}: high-order spectral-element code; widely used for DNS/LES in complex geometries.
\end{itemize}
\end{block}
\begin{itemize}
  \item Choosing a tool depends on: accuracy needs, geometry, turbulence model, and HPC resources.
\end{itemize}
\column{0.44\textwidth}
\includegraphics[width=\textwidth]{pictures/5-1/nek5000-1000x500.jpg}
\end{columns}
\end{frame}

% -------------------------
\begin{frame}{Takeaways}
\begin{itemize}
  \item Turbulence is everywhere, but canonical flows give structure to the chaos.
  \item Each canonical flow highlights different mechanisms (instability, shear, walls, separation, resonance).
  \item Open videos, datasets, and codes make it easier than ever to \emph{see} turbulence and to explore it quantitatively.
\end{itemize}

\vspace{0.8em}

\end{frame}
\end{document}